\chapter{Lua-API-Referenz}
\section{Laufzeitmodell}
app.lua liefert eine Tabelle M. M.loop() ist erforderlich und wird wiederholt aufgerufen; M.init(), M.event(event\_type) und M.cleanup() sind optional. Ohne loop wird die App abgelehnt. Bei Laufzeitfehlern kann die Firmware auf eine Ersatzanwendung wechseln.

Das Sensorskript in sensor/*.lua liefert eine Tabelle sensor mit sensor.get(index), das einen Zahlenwert zurückgibt. Die Firmware stellt sensor global bereit. apppath ist der Pfad zum aktuellen App-Ordner. Module werden mit require("module") geladen; für optionale Module kann pcall(require, "module") verwendet werden.
\section{Ressourcen und Darstellung}
loadimg und loadfont liefern ein Handle oder nil; free gibt geladene Bild- und Schriftressourcen frei. fill füllt den Bildschirm. update überträgt den Frame und löscht anschließend auf Schwarz. drawpolygon benötigt eine gerade Koordinatenanzahl mit mindestens vier Elementen. update einmal pro Frame aufrufen und alle Ressourcen in cleanup freigeben.
\section{Weitere Module}
system steuert Wartezeit, Hintergrundbeleuchtung, V2-Pull-up, 5-V-Ausgang und Schlafkonfiguration und liest Tasten, ADC und Uhrzeit. onewire sucht DS18B20-Sensoren und liest die Temperatur in Celsius. nvs verwendet den NVS-Namensraum app; get liefert eine Zeichenkette oder nil, push einen Wahrheitswert. ble initialisiert BLE und liefert Drehzahl, Geschwindigkeit und Gang.
\section{display}
\begin{itemize}
\item display.loadimg(path) -> img\_handle|nil
\item display.loadfont(path) -> font\_handle|nil
\item display.free(handle)
\item display.fill(color)
\item display.update()
\item display.drawimg(img, x, y)
\item display.string(font, x, y, text, color)
\item display.fillrect(x, y, width, height, color)
\item display.drawline(x1, y1, x2, y2, color)
\item display.drawcircle(x, y, r, color)
\item display.drawfillcircle(x, y, r, color)
\item display.drawroundrect(x1, y1, x2, y2, r, color)
\item display.drawarrow(x0, y0, x1, y1, w, color)
\item display.drawfillarrow(x0, y0, x1, y1, w, color)
\item display.drawpolygon(\{x1, y1, x2, y2, ...\}, color)
\item display.RED
\item display.GREEN
\item display.BLUE
\item display.WHITE
\item display.BLACK
\item display.YELLOW
\item display.CYAN
\item display.PURPLE
\item display.GRAY
\end{itemize}
\section{system}
\begin{itemize}
\item system.delay(ms)
\item system.backlight(level)
\item system.V2\_pullup(enable)
\item system.out5V(enable)
\item system.keys() -> int
\item system.adc\_ch(channel) -> number
\item system.get\_time() -> hour, minute, second
\item system.config\_sleep(enable)
\end{itemize}
\section{onewire}
\begin{itemize}
\item onewire.scan18b20() -> count
\item onewire.read18b20(index) -> temperature\_c
\end{itemize}
\section{nvs}
\begin{itemize}
\item nvs.get(key) -> string|nil
\item nvs.push(key, value\_string) -> boolean
\end{itemize}
\section{ble}
\begin{itemize}
\item ble.init() -> true
\item ble.get\_telemetry() -> rpm, speed, gear
\end{itemize}
\section{Beispiele und Hinweise}
Das Voltmeterbeispiel verwendet Koordinaten innerhalb von 160\,\texttimes\,80 Pixeln und begrenzt die Balkenbreite. Das Sensorbeispiel zeigt ADC-Abfrage und Glättung, keine vollständige NTC-Temperaturberechnung: die Beispieltransformation durch die tatsächliche Kalibrierung ersetzen. sensor.get schnell halten, Kalibrierwerte bei Bedarf über nvs speichern und aus Zeichenketten umwandeln. Pfade aus apppath bilden. DS18B20 einmal suchen und dann im Loop lesen.
\begin{verbatim}
local M = {}
local font = display.loadfont(apppath .. "/../fonts/H24XB.FNT")

local function read_voltage()
return sensor.get(1)
end

function M.loop()
local v = read_voltage()
display.fill(display.BLACK)
display.string(font, 10, 30,
    string.format("%2.1f V", v), display.WHITE)
display.fillrect(10, 60,
    math.max(0, math.min(140, math.floor((v - 9.5) * 12))),
    8, display.GREEN)
display.update()
end

function M.cleanup()
display.free(font)
end

return M
\end{verbatim}
\begin{verbatim}
sensor = {}

local filt = 0.0

function sensor.get(i)
i = i or 0
system.V2_pullup(1)

local raw = system.adc_ch(i)
filt = filt + (raw - filt) / 10.0

local value = filt * 1.0
return value
end

return sensor
\end{verbatim}
